%======================================================================
\section*{Chapter 11: The Analysis of Variance (ANOVA)}
\addcontentsline{toc}{section}{Chapter 11: The Analysis of Variance (ANOVA)}

\subsection*{11.1 Single-Factor ANOVA}
\addcontentsline{toc}{subsection}{11.1 Single-Factor ANOVA}

Single-factor ANOVA focuses on the means of two or more populations. There are $I$ populations (treatments) with means $\mu_1,\mu_2,\dots,\mu_I$:
\[
H_0:\mu_1=\mu_2=\dots=\mu_I,\qquad H_a: \text{at least one } \mu_i \text{ is different from the others}.
\]

\textbf{Notation and assumptions.}
\begin{itemize}
  \item $X_{ij}$: the random variable denoting the $j$th measurement from the $i$th population (treatment).
  \item $x_{ij}$: the observed value of $X_{ij}$.
  \item Different samples are independent of each other, and the $X_{ij}$'s within a sample are also independent of each other.
\end{itemize}
The ANOVA here is developed for the case of equal sample sizes, which is called a \textbf{balanced study}.
Let $J$ be the number of observations in each sample. Then the $i$th sample mean is
\[
\bar{X}_{i\cdot} = \frac{1}{J}\sum_{j=1}^J X_{ij},\qquad i=1,\dots,I,
\]
and, for the $IJ$ observations, the \textbf{grand mean} is defined as
\[
\bar{X}_{\cdot\cdot} = \frac{1}{IJ}\sum_{i=1}^I\sum_{j=1}^J X_{ij} = \frac{1}{IJ}\sum_{i=1}^I J\bar{X}_{i\cdot} = \frac{1}{I}\sum_{i=1}^I\bar{X}_{i\cdot}.
\]

\textbf{ANOVA assumptions.} The $I$ populations (treatments) are normally distributed with the same variance $\sigma^2$.
Then the $X_{ij}$'s are independent and normal with $E(X_{ij})=\mu_i$ and $V(X_{ij})=\sigma^2$.

\textbf{Sums of squares and mean squares.} There are two types of variation:
\begin{enumerate}
  \item between-sample variation,
  \item within-sample variation.
\end{enumerate}

\textbf{Definition.} A measure of between-sample variation is called the \textbf{treatment sum of squares} (SSTr), defined by
\[
\mathrm{SSTr} = \sum_{i=1}^I\sum_{j=1}^J(\bar{X}_{i\cdot}-\bar{X}_{\cdot\cdot})^2 = \sum_{i=1}^I J(\bar{X}_{i\cdot}-\bar{X}_{\cdot\cdot})^2
= J\left[(\bar{X}_{1\cdot}-\bar{X}_{\cdot\cdot})^2+\dots+(\bar{X}_{I\cdot}-\bar{X}_{\cdot\cdot})^2\right].
\]
A measure of within-sample variation is called the \textbf{error sum of squares} (SSE), defined by
\[
\mathrm{SSE} = \sum_{i=1}^I\sum_{j=1}^J(X_{ij}-\bar{X}_{i\cdot})^2 = \sum_{i=1}^I(J-1)S_i^2 = (J-1)\left[S_1^2+\dots+S_I^2\right].
\]

\textbf{Theorem (the ANOVA identity).}
\[
\mathrm{SST} = \mathrm{SSTr}+\mathrm{SSE},\qquad\text{where}\qquad \mathrm{SST}=\sum_{i=1}^I\sum_{j=1}^J(x_{ij}-\bar{x}_{\cdot\cdot})^2
\]
is the \textbf{total sum of squares} ($x_{ij}$ and $\bar{x}_{\cdot\cdot}$ are the observed values).
Here SSTr is the \emph{explained} variation, as it could happen due to a treatment effect; SSE is called the \emph{unexplained} variation,
as it occurs due to randomness within each sample.
If the between-sample variation is large relative to the within-sample variation (more precisely, if MSTr is large relative to MSE, see below), we reject $H_0$.

\textbf{Theorem.} When the ANOVA assumptions are satisfied,
\begin{enumerate}
  \item[(i)] SSE and SSTr are independent;
  \item[(ii)] $\dfrac{\mathrm{SSE}}{\sigma^2}\sim\chi^2_{IJ-I}$;
  \item[(iii)] when $H_0:\mu_1=\dots=\mu_I$ is true, $\dfrac{\mathrm{SSTr}}{\sigma^2}\sim\chi^2_{I-1}$.
\end{enumerate}

\emph{Proof.}
(i) SSTr consists of the $\bar{X}_{i\cdot}$'s and SSE consists of the $S_i^2$'s; under the normality assumption of ANOVA, $\bar{X}_{i\cdot}$ and $S_i^2$ are independent, so SSTr and SSE are independent.

(ii)
\[
\frac{\mathrm{SSE}}{\sigma^2} = \frac{(J-1)\left[S_1^2+\dots+S_I^2\right]}{\sigma^2}
= \underbrace{\frac{(J-1)S_1^2}{\sigma^2}}_{\sim\chi^2_{J-1}}+\dots+\underbrace{\frac{(J-1)S_I^2}{\sigma^2}}_{\sim\chi^2_{J-1}}
\sim\chi^2_{I(J-1)}=\chi^2_{IJ-I},
\]
since the $I$ terms are independent.

(iii) Define $Y_i=\bar{X}_{i\cdot}$. Under $H_0$, $Y_1,\dots,Y_I$ are independent and normally distributed with the same mean and variance $\sigma^2/J$.
So, with $S_Y^2=\frac{1}{I-1}\sum_i(Y_i-\bar{Y})^2$,
\[
\frac{(I-1)S_Y^2}{\sigma^2/J}\sim\chi^2_{I-1},
\qquad
\frac{(I-1)S_Y^2}{\sigma^2/J} = \frac{(I-1)J}{\sigma^2}\cdot\frac{1}{I-1}\sum_{i=1}^I(Y_i-\bar{Y})^2
= \frac{J}{\sigma^2}\sum_{i=1}^I(\bar{X}_{i\cdot}-\bar{X}_{\cdot\cdot})^2 = \frac{\mathrm{SSTr}}{\sigma^2}.
\]
Hence $\mathrm{SSTr}/\sigma^2\sim\chi^2_{I-1}$. \hfill$\square$

Analogous to the definition of the sample variance, we give the following definitions.

\textbf{Definition.} The \textbf{mean square for treatments} (MSTr) and the \textbf{mean square for error} (MSE) are
\[
\mathrm{MSTr}=\frac{\mathrm{SSTr}}{I-1},\qquad \mathrm{MSE}=\frac{\mathrm{SSE}}{IJ-I}.
\]

\textbf{Proposition.} When the ANOVA assumptions are satisfied,
\begin{enumerate}
  \item[(i)] if $H_0$ is true, $E(\mathrm{MSTr})=\sigma^2$, i.e.\ MSTr is then an unbiased estimator of $\sigma^2$ (when $H_0$ is false, $E(\mathrm{MSTr})>\sigma^2$);
  \item[(ii)] $E(\mathrm{MSE})=\sigma^2$, i.e.\ MSE is always an unbiased estimator of $\sigma^2$.
\end{enumerate}

\emph{Proof.} If $X\sim\chi^2_\nu$, then $E(X)=\nu$, since $X=Z_1^2+\dots+Z_\nu^2$ and $E(X)=E(Z_1^2)+\dots+E(Z_\nu^2)=1+\dots+1=\nu$.

(i) $E(\mathrm{MSTr}) = E\!\left(\dfrac{\mathrm{SSTr}}{I-1}\right) = \dfrac{1}{I-1}E(\mathrm{SSTr})$ $\;(**)$.
Under $H_0$ we know $E\!\left(\dfrac{\mathrm{SSTr}}{\sigma^2}\right)=I-1 \Rightarrow E(\mathrm{SSTr})=\sigma^2(I-1)$ $\;(***)$.
Substituting $(***)$ into $(**)$, we get $E(\mathrm{MSTr})=\sigma^2$.

(ii) $E(\mathrm{MSE}) = E\!\left(\dfrac{\mathrm{SSE}}{IJ-I}\right) = \dfrac{1}{IJ-I}E(\mathrm{SSE}) = \dfrac{1}{IJ-I}(IJ-I)\sigma^2 = \sigma^2.$ \hfill$\square$

\textbf{The $F$ test.}
\begin{itemize}
  \item If $H_0$ is true, then MSTr and MSE are about the same, namely $\dfrac{\mathrm{MSTr}}{\mathrm{MSE}}\approx1$.
  \item If $H_a$ is true, then MSTr tends to be larger than MSE, and the ratio $\dfrac{\mathrm{MSTr}}{\mathrm{MSE}}>1$.
\end{itemize}
Recall: if $Y_1\sim\chi^2_{\nu_1}$ and $Y_2\sim\chi^2_{\nu_2}$ are independent, then $F=\dfrac{Y_1/\nu_1}{Y_2/\nu_2}\sim F_{\nu_1,\nu_2}$.

\textbf{Proposition.} When the ANOVA assumptions are satisfied and $H_0:\mu_1=\dots=\mu_I$ is true, the test statistic
$F=\mathrm{MSTr}/\mathrm{MSE}$ has an $F$ distribution with $I-1$ numerator df and $IJ-I$ denominator df.

\emph{Proof.}
\[
F=\frac{\mathrm{MSTr}}{\mathrm{MSE}}=\frac{\mathrm{SSTr}/(I-1)}{\mathrm{SSE}/(IJ-I)}
=\frac{\overbrace{[\mathrm{SSTr}/\sigma^2]}^{\sim\chi^2_{I-1}}/(I-1)}{\underbrace{[\mathrm{SSE}/\sigma^2]}_{\sim\chi^2_{IJ-I}}/(IJ-I)}\sim F_{I-1,\,IJ-I},
\]
using the independence of SSTr and SSE. \hfill$\square$

\textbf{Single-factor ANOVA test.}
\begin{itemize}
  \item $H_0:\mu_1=\mu_2=\dots=\mu_I$.
  \item $H_a$: at least one $\mu_i$ is different from the others.
  \item Test statistic value: $f=\dfrac{\mathrm{MSTr}}{\mathrm{MSE}}=\dfrac{\mathrm{SSTr}/(I-1)}{\mathrm{SSE}/(IJ-I)}$.
  \item Rejection region: $f\ge F_{\alpha,I-1,IJ-I}$.
  \item $P$-value: the area under the $F_{I-1,IJ-I}$ curve to the right of the calculated $f$.
\end{itemize}

\textbf{ANOVA table.}
\begin{center}
\begin{tabular}{lcccc}
\toprule
Source of variation & df & Sum of squares & Mean square & $f$\\
\midrule
Treatments & $I-1$  & SSTr & $\mathrm{MSTr}=\mathrm{SSTr}/(I-1)$ & $\mathrm{MSTr}/\mathrm{MSE}$\\
Error      & $IJ-I$ & SSE  & $\mathrm{MSE}=\mathrm{SSE}/(IJ-I)$ & \\
Total      & $IJ-1$ & $\mathrm{SST}=\mathrm{SSTr}+\mathrm{SSE}$ & & \\
\bottomrule
\end{tabular}
\end{center}

\subsection*{11.2 Multiple Comparisons in ANOVA}
\addcontentsline{toc}{subsection}{11.2 Multiple Comparisons in ANOVA}

When $H_0$ is rejected ($H_a$ is accepted), we may want to find which $\mu_i$'s are different from the others.
A method for such an analysis is called a \textbf{multiple comparisons procedure}. These methods often follow the steps below:
\begin{enumerate}
  \item Calculate a CI for each pairwise difference $\mu_i-\mu_j$ with $i<j$ (there are $\binom{I}{2}$ of them, e.g.\ $\binom{5}{2}=10$ for $I=5$).
  \item If the CI does not include 0, then $\mu_i-\mu_j$ cannot be 0, which means $\mu_i$ and $\mu_j$ are significantly different.
  \item Carry out step 2 for each $\mu_i-\mu_j$ with $i<j$.
\end{enumerate}
We can conduct the preceding steps using different methods; here we introduce Tukey's procedure.

\textbf{Tukey's procedure.}

\textbf{Definition.} Let $Z_1,\dots,Z_m$ be independent standard normal random variables, and let $W$ be independent of them with a chi-squared distribution with $\nu$ df. Then
\[
Q = \frac{\max_{i,j}|Z_i-Z_j|}{\sqrt{W/\nu}} = \frac{\max(Z_1,\dots,Z_m)-\min(Z_1,\dots,Z_m)}{\sqrt{W/\nu}} \tag{1}
\]
has a \textbf{studentized range distribution}. The critical value is denoted $Q_{\alpha,m,\nu}$.

The quantities in this definition are
\[
Z_i=\frac{\bar{X}_{i\cdot}-\mu_i}{\sigma/\sqrt{J}},\qquad m=I,\qquad W=\frac{\mathrm{SSE}}{\sigma^2}=\frac{(IJ-I)\mathrm{MSE}}{\sigma^2},\qquad \nu=IJ-I. \tag{2}
\]
Substituting (2) into (1), we get
\[
Q = \frac{\max_{i,j}\left|\dfrac{\bar{X}_{i\cdot}-\mu_i}{\sigma/\sqrt{J}}-\dfrac{\bar{X}_{j\cdot}-\mu_j}{\sigma/\sqrt{J}}\right|}{\sqrt{\dfrac{(IJ-I)\mathrm{MSE}}{\sigma^2}\Big/(IJ-I)}}
\]
\[
= \frac{\max_{i,j}\left|(\bar{X}_{i\cdot}-\bar{X}_{j\cdot})-(\mu_i-\mu_j)\right|\big/(\sigma/\sqrt{J})}{\frac{1}{\sigma}\sqrt{\mathrm{MSE}}}
= \frac{\max_{i,j}\left|(\bar{X}_{i\cdot}-\bar{X}_{j\cdot})-(\mu_i-\mu_j)\right|}{\sqrt{\mathrm{MSE}/J}}. \tag{3}
\]
We know
\[
P(Q\le Q_{\alpha,I,IJ-I})=1-\alpha. \tag{4}
\]
Now substitute (3) into (4) to obtain the CIs:
\begin{align*}
1-\alpha &= P\!\left(\frac{\max_{i,j}|(\bar{X}_{i\cdot}-\bar{X}_{j\cdot})-(\mu_i-\mu_j)|}{\sqrt{\mathrm{MSE}/J}}\le Q_{\alpha,I,IJ-I}\right)\\
&= P\!\left(\frac{|(\bar{X}_{i\cdot}-\bar{X}_{j\cdot})-(\mu_i-\mu_j)|}{\sqrt{\mathrm{MSE}/J}}\le Q_{\alpha,I,IJ-I}\ \text{ for all } i,j\right)\\
&= P\!\left(-Q_{\alpha,I,IJ-I}\le\frac{(\bar{X}_{i\cdot}-\bar{X}_{j\cdot})-(\mu_i-\mu_j)}{\sqrt{\mathrm{MSE}/J}}\le Q_{\alpha,I,IJ-I}\ \text{ for all } i,j\right)\\
&= P\Big(\bar{X}_{i\cdot}-\bar{X}_{j\cdot}-Q_{\alpha,I,IJ-I}\sqrt{\mathrm{MSE}/J}\le\mu_i-\mu_j\\
&\qquad\quad \le\bar{X}_{i\cdot}-\bar{X}_{j\cdot}+Q_{\alpha,I,IJ-I}\sqrt{\mathrm{MSE}/J}\ \text{ for all } i,j\Big).
\end{align*}

\textbf{Proposition.} For each $i<j$, form the interval
\[
\bar{x}_{i\cdot}-\bar{x}_{j\cdot}\pm Q_{\alpha,I,IJ-I}\sqrt{\mathrm{MSE}/J}.
\]
There are $\binom{I}{2}$ such intervals, and their \textbf{simultaneous} confidence level is $100(1-\alpha)\%$.
The margin of error $Q_{\alpha,I,IJ-I}\sqrt{\mathrm{MSE}/J}$ is called \textbf{Tukey's honestly significant difference} (HSD).

\textbf{Tukey's method for identifying significantly different $\mu_i$'s.}
\begin{enumerate}
  \item List the sample means in increasing order, together with the corresponding population (treatment).
  \item Starting from the leftmost (smallest) mean, compare the other means with it. If the CI for a difference includes 0, underscore those means together (they are not significantly different).
  \item Do the same for the other means. Then we can judge which $\mu_i$'s are significantly different.
\end{enumerate}

\textbf{Example.} Suppose that $I=5$ with $\bar{x}_{2\cdot}\le\bar{x}_{5\cdot}\le\bar{x}_{1\cdot}\le\bar{x}_{3\cdot}\le\bar{x}_{4\cdot}$. Let us say
\begin{itemize}
  \item the CI for $\mu_2-\mu_5$ includes 0,
  \item the CI for $\mu_1-\mu_2$ includes 0,
  \item the CI for $\mu_2-\mu_3$ does not include 0,
  \item the CI for $\mu_2-\mu_4$ does not include 0,
  \item the CI for $\mu_1-\mu_3$ includes 0.
\end{itemize}
Then, in the order of the sample means,
\begin{center}
\begin{tabular}{lccccc}
\toprule
Group & 2 & 5 & 1 & 3 & 4\\
Sample mean & $\bar{x}_{2\cdot}$ & $\bar{x}_{5\cdot}$ & $\bar{x}_{1\cdot}$ & $\bar{x}_{3\cdot}$ & $\bar{x}_{4\cdot}$\\
\bottomrule
\end{tabular}
\end{center}
treatments 2, 5 and 1 are underscored together (not significantly different), and $\mu_2$ is significantly different from $\mu_3$ and $\mu_4$.

For the second step in Tukey's method, instead of constructing the CI for each difference, we can check the condition
\[
|\bar{x}_{i\cdot}-\bar{x}_{j\cdot}| < Q_{\alpha,I,IJ-I}\sqrt{\mathrm{MSE}/J}.
\]
If this is true, then the CI includes 0.

\textbf{CIs for other parametric functions.} Consider $\theta=\sum_i c_i\mu_i$, where the $c_i$'s are constants.
The natural estimator of $\theta$ is $\hat\theta=\sum_i c_i\bar{X}_{i\cdot}$. We need $\dfrac{\hat\theta-\theta}{\sigma_{\hat\theta}}$, so we first find $V(\hat\theta)$:
\[
V(\hat\theta) = V\!\left(\sum_i c_i\bar{X}_{i\cdot}\right) = \sum_i c_i^2V(\bar{X}_{i\cdot}) = \frac{\sigma^2}{J}\sum_i c_i^2,
\qquad
\sigma_{\hat\theta}=\sqrt{\frac{\sigma^2}{J}\sum_i c_i^2}.
\]
So we have $\dfrac{\sum_i c_i\bar{X}_{i\cdot}-\sum_i c_i\mu_i}{\sqrt{\frac{\sigma^2}{J}\sum_i c_i^2}}$.
Since $\sigma$ is unknown and its estimator is $\sqrt{\mathrm{MSE}}$ (MSE estimates $\sigma^2$), we write the CI
\[
\sum_i c_i\bar{x}_{i\cdot}\pm t_{\alpha/2,\nu}\sqrt{\frac{\mathrm{MSE}}{J}\sum_i c_i^2},\qquad \nu=IJ-I.
\]
